% ###################################################################################################
% Chapter 9 - Game Master Tools
% ###################################################################################################
\chapter{Game Master Tools}

Running a game of Eidos requires a shift in perspective for Game Masters (GMs) accustomed to highly structured, grid-locked tabletop systems. In Eidos, the story and mechanics are inseparably bound. The GM is not an adversary but a weaver of consequence, translating player rolls and decisions into a living, responsive world. 

This chapter outlines essential guidance, math-scaling principles, and quick-generation tools to help you run compelling campaigns.

\section{Worldbuilding in a Narrative RPG}

Eidos is designed by default to simulate a pre-modern, low-fantasy environment. In such settings, magic is rare and terrifying, the wilderness is deadly, and local cultures and religious institutions hold immense social power.

\subsection{Start Small, Build Deep}
Instead of mapping an entire continent, focus on a single, localized conflict.
\begin{itemize}
    \item \textbf{The Microcosm:} Create one village or town nestled in a dangerous region. Map out its primary power structures: who runs the church? Who controls the military guard? What local cults lurk in the shadows?
    \item \textbf{Theme over Scale:} Emphasize sensory details—the smell of wet horse hair, the mud clinging to leather boots, the distant ringing of temple bells. This grounds the survival mechanics of Chapter 8 in physical reality.
\end{itemize}

\subsection{Leveraging Character Sheets}
GMs should actively look at the characters' **Alignment** (Chapter 1) and **Spirituality** (Chapter 2).
\begin{itemize}
    \item If a character is an *Orthodox Adherent* of a *Cult*, introduce a rival religious sect that treats them with suspicion.
    \item Use a character's *Personality Quirks* to create interesting social challenges. A character with a fear of heights will find a mountain pass far more intimidating, prompting narrative actions.
\end{itemize}

\section{NPC and Adversary Design}

Adversaries in Eidos do not need complex, multi-page character sheets. They are grouped into three distinct categories based on their role in the story:

\subsection{Minions}
Minions represent common rabble, weak guards, small wild beasts, or minor cultists. They exist to form crowds or test the players' basic tactical capabilities.
\begin{itemize}
    \item \textbf{Stats:} Static Defense of 10. Attack modifier of +0.
    \item \textbf{1-Hit KO:} Minions do not track wounds or roll on injury tables. Any successful physical attack that hits them immediately neutralizes them (death, unconsciousness, or routing, depending on the narrative context).
\end{itemize}

\subsection{Elites}
Elites are seasoned soldiers, trained assassin scouts, or ferocious beasts. They pose a significant threat to a lone character.
\begin{itemize}
    \item \textbf{Stats:} Passive Defense of $12 + \text{AGI}$. Attack check modifier of +2 to +4. 
    \item \textbf{Attributes:} Standard Array (e.g., AGI +2, STR +1).
    \item \textbf{Wounds:} Elites roll on Critical Injury tables normally. They are neutralized when they suffer an "Instant Death" outcome or when their physical attributes are reduced to -4 due to severe trauma.
\end{itemize}

\subsection{Bosses}
Bosses are legendary figures—celebrated warlords, high inquisitors, or ancient monsters. 
\begin{itemize}
    \item \textbf{Stats:} Passive Defense of $14 + \text{AGI}$. Attack check modifier of +6 or higher.
    \item \textbf{Resolve Points (3):} Bosses possess 3 Resolve Points. They can spend 1 Resolve Point to completely shrug off a debilitating status effect or downgrade a critical injury roll (e.g., converting a severed arm or death roll into a deep cutting chest slash that takes 1d4 weeks to heal).
    \item \textbf{Legendary Reactions:} A Boss can take a reaction after any player's turn, allowing them to reposition, shout an intimidating command, or make a quick strike.
\end{itemize}

\section{Encounter Scaling and Difficulty}

Eidos does not use strict "Challenge Ratings." Instead, GMs balance encounters by adjusting Target Numbers (TNs) and the composition of adversaries.

\subsection{Combat Encounters}
To design a balanced combat encounter, use the following composition guidelines based on a standard party of 4 characters:
\begin{itemize}
    \item \textbf{Easy Encounter:} 4–6 Minions. Ideal for introducing players to combat flow.
    \item \textbf{Moderate Encounter:} 2 Elites and 2–3 Minions. A solid challenge requiring teamwork.
    \item \textbf{Hard Encounter:} 1 Boss and 2 Elites. Extremely lethal; players must utilize active defense, shields, and tactical planning.
\end{itemize}

\subsection{Social and Environmental Encounters}
When setting target numbers for non-combat actions (e.g., climbing a cliff, deciphering an ancient text, or bribing a guard):
\begin{itemize}
    \item **Moderate Tasks (TN 15):** The standard benchmark. An untrained character has a 50\% chance of success; a trained character succeeds regularly.
    \item **Hard Tasks (TN 20 or 25):** Requires specialized skills. Players must seek modifiers (e.g., ropes for climbing, or a diplomatic gift for bribing) to increase their chances.
\end{itemize}
